\section{Conclusions and Future Work}

\subsection{Conclusions}
This research investigated to which extent the storage and analysis efficiency can be improved by transforming a daily
historical DNS record database into a delta-based storage system. The results confirmed the underlying hypothesis that 
a majority of the collected DNS records remain unchanged between measurements. The introduced DeltaParq, a schema-aware 
delta-based storage prototype build for the OpenINTEL project, showed that this can be used to reduce storage usage.
An evaluation over multiple datasets showed that up to about 49.25\% of storage space can be saved using the 
Progressive strategy, which also comes with a significantly longer analysis time with an increase of 98.65\%. Other evaluated strategies, 
such as the BST strategy, show that saving a substantial amount of storage space can be done with less impact 
on the analysis time. Using this the system can be tuned to the preferences of its user.

% Hypothesis deemed correct
% Succesfully constructed the DeltaParq prototype
% Maximum gain lies around 50% but comes with a loss
% However storage gain definetly possible with less removal from analysis time (BST)

\subsection{Future Work}

The findings of this research point to multiple directions for future research. Firstly, the bottleneck identified in Section~\ref{sec:performance_constraints}, 
the shuffle overhead caused by PySpark not being able to use sorted or partitioned data without a Hive metastore. A deployment with support for partitioning on storing and reading may significantly 
reduce the overhead in analysis time, which could potentially make the analysis time of strategies that use long dependency chains, such as 
the Progressive strategy, significantly lower. Furthermore, it may make more frequent anchor resets, such as weekly instead of monthly, 
more viable as the increased analysis time identified in~\ref{sec:discussion} would have less of an impact.

Secondly, this research evaluated the DeltaParq prototype on only 2 datasets based on a single underlying structure. Evaluating the 
schema-aware delta-based storage systems over more datasets, both within the same OpenINTEL structure, and other domains with similar 
volatile and non-volatile data columns, would strengthen the conclusions and the generalizability beyond the OpenINTEL system.

Finally, strategies can be optimized for both better performance and data safety.
Further research could explore the possibility of using the structure of the underlying dataset combined with the 
properties of the strategies to automatically select a strategy that performs optimally for a user defined balance 
between storage reduction and analysis time increase. Furthermore, certain strategies could be extended. For example,
adding a bridging delta between the last day of the month and the first day of the next month in the progressive strategy, 
would allow for a smart system to recover a singular corrupted delta.


% Spark optimization -> May make more aggressive strategies possible
% Test on more datasets
% Create a model dataset -> Strategy
